\chapter{Conclusion}\label{ch: Conclusion}
This report focused on simulating the basic BTW model, exploring its steady-state behaviour and the frequency distributions of four avalanche properties, and examining its Abelian property. An extended model with gradient-dependent toppling rules was also simulated, and its steady-state behaviour and avalanche frequency distributions were compared with those of the basic model.

For the basic BTW model, the average height fluctuates around 2.1 after about 10,000 grain additions. The distribution of Topples exhibits approximate power-law behaviour. The distribution of positive Area values is reasonably approximated by a log-normal distribution, while positive Loss values are approximately described by a geometric distribution. A log-normal distribution provides only a rough approximation for positive Length values. Spearman correlations were also calculated between the four properties. Topples, Area and Length are strongly positively correlated with one another, with coefficients above 0.9. In contrast, Loss has weaker positive correlations with these properties, with coefficients of approximately 0.63.

For the extended model, the average height fluctuates around 26.5 after about 190,000 grain additions. The spatial distribution of sand forms a pyramid-shaped pile, with heights generally increasing towards the centre. The frequencies of Topples and Area decrease approximately linearly up to about 25, followed by a tail of low-frequency events. Length also shows an approximately linear decrease in frequency, without a similar tail in the observed range. The observed Loss values are limited to 0, 1 and 2, with Loss = 1 accounting for 96.22\% of the 10,000 observations after burn-in.

Overall, the comparison shows that changing the local toppling rules can largely affect both the spatial distribution of sand and avalanche statistics, even when the average height in both models eventually stabilises.